%%%%%%%%%%%%%%%%%%%%%%%%%%%%%%%%%%%%%%%%%%%%%%%%%%%%%%%%%%%%%%%%%%%%%%
% Lớp 8 | Chương II. Hằng đẳng thức đáng nhớ và ứng dụng | Hiệu hai bình phương. Bình phương của một tổng hay một hiệu
% Dạng 8D2?6-3: Tính nhanh, tính giá trị biểu thức
%%%%%%%%%%%%%%%%%%%%%%%%%%%%%%%%%%%%%%%%%%%%%%%%%%%%%%%%%%%%%%%%%%%%%%


% nguồn: dan_tu_AI (nhập 14/11/2026)
\begin{bt}%[8D2H6-3]%[ma:cfeaff1c]
	Tính nhanh:\\
	a) $101^2$;\qquad b) $99^2$;\qquad c) $52\cdot 48$.
	\loigiai{
		a) $101^2=(100+1)^2=100^2+2\cdot 100\cdot 1+1^2=10\,000+200+1=10\,201$.\\
		b) $99^2=(100-1)^2=100^2-2\cdot 100\cdot 1+1^2=10\,000-200+1=9\,801$.\\
		c) $52\cdot 48=(50+2)(50-2)=50^2-2^2=2\,500-4=2\,496$.
	}
\end{bt}

% nguồn: dan_tu_AI (nhập 14/11/2026)
\begin{bt}%[8D2H6-3]%[ma:2055c9b7]
	Tính nhanh:\\
	a) $199^2$;\qquad b) $103\cdot 97$;\qquad c) $45^2+10\cdot 45+25$.
	\loigiai{
		a) $199^2=(200-1)^2=40\,000-400+1=39\,601$.\\
		b) $103\cdot 97=(100+3)(100-3)=100^2-3^2=10\,000-9=9\,991$.\\
		c) $45^2+10\cdot 45+25=45^2+2\cdot 45\cdot 5+5^2=(45+5)^2=50^2=2\,500$.
	}
\end{bt}

% nguồn: dan_tu_AI (nhập 14/11/2026)
\begin{bt}%[8D2H6-3]%[ma:3dc50788]
	Tính giá trị của các biểu thức sau:\\
	a) $A=x^2-6x+9$ tại $x=103$;\\
	b) $B=4x^2+4x+1$ tại $x=24{,}5$;\\
	c) $C=x^2-y^2$ tại $x=87$, $y=13$.
	\loigiai{
		a) $A=(x-3)^2$. Tại $x=103$: $A=(103-3)^2=100^2=10\,000$.\\
		b) $B=(2x+1)^2$. Tại $x=24{,}5$: $B=(2\cdot 24{,}5+1)^2=50^2=2\,500$.\\
		c) $C=(x-y)(x+y)$. Tại $x=87$, $y=13$: $C=(87-13)(87+13)=74\cdot 100=7\,400$.
	}
\end{bt}

% nguồn: dan_tu_AI (nhập 14/11/2026)
\begin{bt}%[8D2V6-3]%[ma:8b71c575]
	Cho hai số $x$, $y$ thỏa mãn $x+y=7$ và $xy=10$. Tính:\\
	a) $x^2+y^2$;\qquad b) $(x-y)^2$;\qquad c) $x^2-y^2$, biết thêm $x>y$.
	\loigiai{
		a) $x^2+y^2=(x+y)^2-2xy=7^2-2\cdot 10=29$.\\
		b) $(x-y)^2=x^2+y^2-2xy=29-20=9$.\\
		c) Từ $(x-y)^2=9$ và $x>y$ suy ra $x-y=3$.\\
		Do đó $x^2-y^2=(x-y)(x+y)=3\cdot 7=21$.
	}
\end{bt}

% nguồn: dan_tu_AI (nhập 14/11/2026)
\begin{bt}%[8D2V6-3]%[ma:4dfd8aa4]
	Một sân chơi hình vuông có cạnh $98\text{ m}$. Người ta sửa lại thành sân hình chữ nhật có chiều dài $100\text{ m}$ và chiều rộng $96\text{ m}$.\\
	a) Dùng hằng đẳng thức, tính nhanh diện tích sân chơi lúc đầu.\\
	b) Dùng hằng đẳng thức, tính nhanh diện tích sân chơi sau khi sửa. Diện tích sân chơi thay đổi thế nào?
	\loigiai{
		a) Diện tích lúc đầu:
		\[ 98^2=(100-2)^2=10\,000-400+4=9\,604\ \left(\text{m}^2\right). \]
		b) Diện tích sau khi sửa:
		\[ 100\cdot 96=(98+2)(98-2)=98^2-2^2=9\,604-4=9\,600\ \left(\text{m}^2\right). \]
		Vậy diện tích sân chơi giảm đi $9\,604-9\,600=4\ \left(\text{m}^2\right)$.
	}
\end{bt}

% nguồn: dan_tu_AI (nhập 16/11/2026)
\begin{bt}%[8D2V6-3]%[ma:467340e2]
	Tính diện tích mặt cắt của một ống nước hình vành khuyên, biết bán kính ngoài là $R = 55\text{ mm}$ và bán kính trong là $r = 45\text{ mm}$. (Lấy $\pi \approx 3{,}14$).
	\loigiai{
		Diện tích mặt cắt của ống nước (hình vành khuyên) được tính bởi công thức:
		\[ S = \pi R^2 - \pi r^2 = \pi (R^2 - r^2). \]
		Thay số vào và áp dụng hằng đẳng thức hiệu hai bình phương để tính nhanh:
		\begin{align*} 
			S &\approx 3{,}14 \cdot (55^2 - 45^2) \\
			&= 3{,}14 \cdot (55 - 45)(55 + 45) \\
			&= 3{,}14 \cdot 10 \cdot 100 \\
			&= 3{,}14 \cdot 1\,000 \\
			&= 3\,140 \text{ (mm}^2\text{)}.
		\end{align*}
		Vậy diện tích mặt cắt của ống nước khoảng $3\,140\text{ mm}^2$.
	}
\end{bt}
